\vspace{-0.07in}
\section{A Solvable Case with Zipf-distributed Data}
%Power-Law Spectrum Structure}
\label{sec:one_hot}
% \haodong{We should show the rationale why we choose the sovable model, \citet{kunstner2025scaling} for example here is a zipf law theory paper}
% In this section, we focus on a solvable case where the linear regression problem exhibits a power-law spectrum structure.
% In addition to the strongly‐convex case, here we study a non‐strongly‐convex linear regression model under a power‐law spectrum assumption.  This setup more closely mirrors practical large‐scale deep learning training and is also common in theoretical analyses of scaling laws~\citep{canatar2021spectral,spigler2020asymptotic,maloney2022solvable,cui2021generalization,brandfonbrener2024loss,lin2024scaling}.  Unlike the usual continuous power‐law spectrum, we assume a one‐hot discrete data distribution, which renders the data permutations effectively exchangeable.  As a result, taking expectations becomes tractable and we avoid the challenging random‐matrix concentration arguments required in \Cref{sec:strongly_convex}. We now introduce the precise problem setup.

Natural data distributions often exhibit power law structures. To capture this phenomenon, we go beyond the strongly convex case and analyze a stylized linear regression model with Zipf-distributed data, where the excess risk admits a closed-form expression and the effective reuse rate can be characterized explicitly.

Through this setup, we can see that the effective reuse rate exhibits a similar scaling behavior: 
as the number of epochs $K$ increases, 
$E(K,N)$ initially grows linearly but eventually saturates at a problem-dependent value that increases with $N$. In contrast to the strongly convex case, however, the saturation point does not scale as $\sim \log N$ but instead scales as a power of $N$.

\paragraph{Problem Setup.}
%We extend the strongly convex linear regression setting in \Cref{sec:setup}. 
We use the same notation for excess risk, one-pass and multi-epoch SGD, and $i.i.d.$ training data as in \Cref{sec:setup}.
We specify the data distribution as a Zipf distribution over $d$ one-hot data points, where the $i$-th data point is $\vx^{(i)} = \mu_i \ve_i$
for some $\mu_i > 0$
and the probability of sampling the $i$-th data point is $p_i = c \cdot i^{-\alpha}$ for some constants $c>0$ and $\alpha>1$.
% \kaifengclaw{Notation consistency: later Assumption~\ref{asp:prior-probability-and-norm} uses parameters $a,b$ instead of $\alpha$; consider adding one line mapping between these symbols.}
The label is generated by $y = \inp{\vw^*}{\vx}$ with no label noise.
The ground-truth weight $\vw^*\in\R^d$ follows an isotropic prior distribution.
\begin{assumption}[Parameter Prior] \label{asp:parameter-prior}
    $\vw^*$ is sampled from a prior distribution with 
    %We assume that
    $\E[\vw^*\vw^{*\top}] = \mI$.
\end{assumption}
% We define the optimal expected excess risk as the minimum of the expected excess risk over all positive learning rates:
% \[
% \bar\cR^*(K,N):=\min_{\eta\geq 0}\,\bar\cR(K,N;\eta).
% \]
% Following the same definition as in \Cref{sec:setup},
% the effective reuse rate is defined as
% $E(K,N) := \frac{1}{N}\min\{N' \ge 0: \bar{\cR}^*(1, N') \leq\bar{\cR}^*(K, N)\}$.

\paragraph{Interpretation.} 
%In the above setup, $\mu_i$ can be seen as a reweighting factor for the $i$-th data point.
This setup can be interpreted as a simplified model of real-world data with heavy-tailed feature distributions. Each coordinate represents an atomic feature that appears with Zipf-distributed probability, mimicking the long-tailed statistics observed in domains such as text and natural language. The scaling factors $\mu_i$ encode feature importance, which may reflect, for instance, effects introduced by feature weighting or normalization.



\begin{comment}
-------------------
In this section, we extend our results in strongly convex linear regression to a non–strongly–convex linear regression with a discrete one‑hot power‑law spectrum, similar to the previous theoretical analyses~\citep{hutter2021learning}. Compared to the previous power-law analysis~\citep{canatar2021spectral,spigler2020asymptotic,maloney2022solvable,cui2021generalization,brandfonbrener2024loss,lin2024scaling}, we consider a more concrete analysis. Specifically, we consider a data distribution appears in text data called \textit{Zipf's Law}, where the frequency of the $r$-th most frequent word $x_r$ is proportional to $\frac{1}{r^b}$ and the exponent $b$ is some positive constant~\citep{piantadosi2014zipf,kunstner2024heavy,kunstner2025scaling}.
% This assumption makes data‑shuffle permutations exchangeable, so expectations admit closed‑form expressions and we avoid the random‑matrix concentration in \Cref{sec:strongly_convex}.
%whd: I dont understand the above argument, starting from "this assumption" since no assumption is mentioned here.
The precise setup is described below.

%\subsection{Problem Setup.} 
% \paragraph{Problem Setup.} We extend the strongly convex linear regression setting in Section~\ref{sec:setup}. Here we share the same excess risk, one-pass SGD, multi-epoch SGD training algorithm, and the $i.i.d.$ training data generation notations, but we assume that the data input $\vx\in \R^d$ follows a one-hot distribution. Specifically, let $\{\ve_i\}_{i=1}^n$ denote the standard orthonormal basis of $\mathbb{R}^n$, where $\ve_i^\top \ve_j = \delta_{ij}$. We assume that the input data $\vx$ follows a discrete distribution with $\vx = \mu_i \ve_i$ occurring with probability $p_i$ for $i \in \{1, \dots, d\}$. Consequently, the Hessian matrix takes the form $\mH = \sum_{i=1}^d \mu_i^2 p_i \ve_i \ve_i^{\top} = \text{diag}(\mu_i^2 p_i)_{i=1}^d$. We further define the data matrix $\mLambda := \diag(\mu_i^2)_{i=1}^d$ and the probability matrix $\mP := \diag(p_i)_{i=1}^d$ to describe the data input distribution. Throughout this section, we assume no label noise ($\sigma=0$), and $\mH_{ii} = \frac{\Tilde{c}}{i^a}$ for two constants $\Tilde{c}>0$ and $a>1$ for $i \in \{1, \dots, d\}$. It is worth noting that, in this section, the definition of optimal expected excess risk can be generalized for a broader range of learning rates as
% \begin{align*}
%     \bar\cR^*(K,N):= \min_{\eta \in \R^+} \{\bar\cR(K,N;\eta)\}.
% \end{align*}

\paragraph{Problem Setup.}
We extend the strongly convex linear regression setting in \Cref{sec:setup}. 
We keep the same notation for excess risk, one-pass and multi-epoch SGD, and $i.i.d.$ training data. 
We assume the input $\vx\in\R^d$ is one-hot, letting $\{\ve_i\}_{i=1}^d$ be the standard basis, $\vx=\mu_i\ve_i$ with probability $p_i$ for $i=1,\dots,d$. 
Then the Hessian $\mH=\sum_{i=1}^d \mu_i^2 p_i\,\ve_i\ve_i^\top=\diag(\mu_i^2 p_i)=\mLambda\mP,$
where $\mLambda:=\diag(\mu_i^2)$ and $\mP:=\diag(p_i)$. 
We assume no label noise ($\sigma=0$) and a power-law spectrum $\lambda_i=\tilde c\, i^{-a}$ with constants $\tilde c>0$ and $a>1$. The definition of optimal expected excess risk can be generalized for a broader range of learning rates as
\[
\bar\cR^*(K,N):=\min_{\eta>0}\,\bar\cR(K,N;\eta).
\]
Also, the ground-truth weight $\vw^*\in\R^d$ follows a certain prior distribution as stated in the following assumption.
\begin{assumption}[Parameter Prior] \label{asp:parameter-prior}
    The ground truth $\vw^*$ is sampled from a prior distribution with 
    %We assume that
    $\E[\vw^*\vw^{*\top}] = \mI$.
\end{assumption}

%Distinct from \Cref{sec:setup}, now we consider the regime where the dimension $d$ is asymptotically large; otherwise, the scenario can be approximately considered as a strongly-convex case, which is formally stated below. 
% \begin{assumption}[Large Dimension] \label{asp:large-dimension}
%     We assume that the dimension $d$ is sufficiently large such that $d = \Omega(T^{\frac{1}{a}})$.
% \end{assumption}

% \tingkai{Delete this assumption.}

% Moreover, we make the following assumptions on the distribution of $\vw^*\in\R^d$, which is defined as the ground truth in \Cref{sec:setup}. 
% % \tingkai{cite \citep{lin2024scaling}}
% \begin{assumption}[Parameter Prior] \label{asp:parameter-prior}
%     We assume that $\vw_0 = \boldsymbol{0}$, and $\E[\vw^*\vw^{*\top}] = \mI$.
% \end{assumption}
% Another crucial assumption is that the dimension $d$ needs to be sufficiently large, otherwise the scenario can be approximately considered as a strongly-convex case.
% \begin{assumption}[Large Dimension] \label{asp:large-dimension}
%     We assume that the dimension $d$ is sufficiently large such that $d = \Omega(T^{\frac{1}{a}})$.
% \end{assumption}

%Similar to the proof insights in \Cref{sec:strongly_convex}, the first move to get the formula of the effective reuse rate is to get an accurate proxy of the excess risk. Here,  leveraging the simplicity of the setting, we can derive a general closed formula for the excess risk.

% \begin{lemma} \label{thm:one-hot-loss-estimate}
%     Under \Cref{asp:parameter-prior}, the excesss risk for $K$-epoch training over $N$ fresh data , with learning rate $\eta$ can be given by
%     \begin{align*}
%         \E[\cR(\vw_{K, N})] = \frac{1}{2}\inp{\mP\Lambda}{\left(\mI-\mP+\mP\left(\mI-\eta\Lambda\right)^{2K}\right)^N},
%     \end{align*}
%     where the expectation is over the randomness of $\vw^*$ and training datasets $\{\vx_i,y_i\}_{i=0}^{N-1}$.
% \end{lemma}
% The above lemma states that we can explicitly write out the exact expression for the excess risk. From the above expression for the excess risk, we can observe that, in the absence of label noise interference, and under the condition that the absolute values of all elements of the diagonal matrix $\mI-\eta\Lambda$ are less than $1$, the optimal learning rate can be of the constant order. Therefore, in the subsequent study of the effective reuse rate, we consider using the same learning rate $\eta = \Theta(1)$ for both multi-epoch and one-pass SGD.
% % In this setting, we employ a constant learning rate $\eta = \Theta(1)$ for both the multi-epoch and one-pass case since the system is noiseless. However, $\eta$ cannot be chosen arbitrarily large due to stability constraints. We also denote $E_{\eta, \eta}(K, N)$ as $E(K, N)$. \haodong{Maybe we should not say that we denote $E_{\eta,\eta}$ as $E$, too confusing and maybe not well define}

% It is worth noting that here we are actually describing a more general problem setting than the Zipf law, as we only impose constraints on the power spectrum of the Hessian matrix $\mH$. In contrast, the probability matrix $\mP$ can follow Zipf's law or any other law. In the remainder of this section, we first consider the classic Zipf's law setting, where $\mP$ follows a power law, and the data matrix $\Lambda$ also follows a power law, which is consistent with the previous power law analysis. In \Cref{sec:log-power-spectrum}, we explore the case where $\mP$ follows a log-power spectrum~\citep{lin2024scaling}, and investigate the impact of changing the spectrum's properties on the resulting effective reuse rate formula.
% %\paragraph{Key Proof Insight.} Notice that in the one-hot case, the matrices are commutative and we are able to calculate the high-order expectations of $\vx$ directly. The proof is given in ...

%\subsection{Derivation of the Effective Reuse Rate Under Different Spectrums}

%\subsubsection{Power-Law Spectrum}

%We then provide the characterization of the effective reuse rate under a power-lawspectrum assumption for $\Lambda$ and $\mP$~\citep{kunstner2025scaling}.
    
\end{comment}
%We then present our main results for the effective reuse rate under specific data distributions.  
\subsection{Results on Power-Law Spectrum}\label{sec:power-spectrum}
\begin{assumption}[Power-Law Spectrum] \label{asp:prior-probability-and-norm}
    There exist two constants $a,b> 0$ with $a-b>1$ such that the data input distribution satisfies that $p_i = c i^{-(a-b)}$ and $\Lambda_i = i^{-b}$, where $c = \left(\sum_{i=1}^d \frac{1}{i^{a-b}}\right)^{-1}$.
\end{assumption}
%Under the above assumptions, we can get a characterization of population excesss risk and $E_K$.
%\subsubsection{Theorem of $E(K,N)$}
% \haodong{delete the following part?}
Here we establish matching upper and lower bounds for $\Bar{\cR}^*(K,N)$ in the small-$K$ and large-$K$ regimes, given the solvable model. Comparing with the strongly convex case, we observe a different scaling behavior: when $K \ll N^{\tfrac{b}{a-b}}$, $\Bar{\cR}^*(K,N)$ decays as a power law in $KN$, with exponent $\tfrac{a-1}{a}$;
whereas when $K \gg N^{\tfrac{b}{a-b}}$, $\Bar{\cR}^*(K,N)$ exhibits a power-law decay in $N$ and is independent of $K$.
\begin{theorem} \label{thm:one-hot-scaling-law}
    Consider a $K$-epoch SGD over $N$ fresh data. Under Assumptions~\ref{asp:parameter-prior}-\ref{asp:prior-probability-and-norm}, and given the data dimension $d = \Omega((KN)^{\frac{1}{a}})$, it holds that
    % \kaifengclaw{Math-check: please ensure the same dimension condition is also sufficient for the one-pass comparator used in $E(K,N)$; otherwise the inversion from risk to $E(K,N)$ may need an extra assumption statement.}
    \[
    \Bar{\cR}^*(K,N) \asymp 
    \begin{cases}
         \left(KN\right)^{-\frac{a-1}{a}} & \text{for } K = o(N^{\frac{b}{a-b}}) \\
         N^{-\frac{a-1}{a-b}} & \text{for } K = \omega(N^{\frac{b}{a-b}}).
    \end{cases}
    \]
\end{theorem}

% \haodong{Some interpretations about the results,  the exponents $a,b$}
Then we derive the formula of $E(K, N)$ by first solving the equation $\Bar{\cR}^*(1,T')=\Bar{\cR}^*(K,N)$ based on \Cref{thm:one-hot-scaling-law}, and divide $T'$ by $N$.
\begin{theorem}[Multi-Epoch Scaling Under Power-Law Spectrum] \label{thm:one-hot-E_K-v2}
    Consider a $K$-epoch SGD over $N$ fresh data. Under Assumptions~\ref{asp:parameter-prior}-\ref{asp:prior-probability-and-norm}, and given the data dimension $d = \Omega((KN)^{\frac{1}{a}})$, it holds that
    \[
    E(K, N) =
    \begin{cases}
         K(1 + o(1)) & \text{for } K = o(N^{\frac{b}{a-b}}) \\
         \Theta(N^{\frac{b}{a-b}}) & \text{for } K = \omega(N^{\frac{b}{a-b}}).
    \end{cases}
    \]
    % \kaifengclaw{Math-check: in the large-$K$ regime this is written as $\Theta(N^{\frac{b}{a-b}})$; if possible, add a brief derivation line showing how the exponent comes from equating the two risk scalings, to avoid exponent slips.}
\end{theorem}
Under the assumption of a logarithmic power-law spectrum, the trend of the effective reuse rate as a function of \( K \) approximates the phenomena described in \Cref{thm:convex-E_K} in the strongly convex setting and the trend described in \Cref{thm:one-hot-E_K-v2} under the power-law spectrum assumption.
% \kaifengclaw{Logic/reference typo: the second citation likely should point to Theorem~\ref{thm:one-hot-E_K-v2}, not again to \Cref{thm:convex-E_K}.} 
We still observe an effective-reuse regime $(E(K, N) \approx K)$ when \( K \) is relatively small $( K \ll N^{b/(a-b)})$, and as \( K \) increases, the effective reuse rate undergoes a phase transition, converging to an upper bound determined by \( N \), entering the limited-reuse regime $(E(K, N) =  \Theta(N^{b/(a-b)}))$. 

We can see that the exponent of this power of $N$ is determined by the rate of eigenvalue decay of the Hessian and the rate of norm decay of the parameter with respect to dimension. The proofs of \Cref{thm:one-hot-scaling-law} and \Cref{thm:one-hot-E_K-v2} are given in Appendix~\ref{sec:proof-one-hot-scaling-law-power-spectrum} and Appendix~\ref{sec:proof-one-hot-Ek-power-spectrum} respectively.

% \Cref{thm:one-hot-E_K-v2} suggests that, under a power-law spectrum structure, the effective reuse rate exhibits a scaling behavior similar to that of the strongly convex case. Specifically, when the number of epochs is smaller than some power of the fresh data size $(K \ll N^{\frac{b}{a-b}})$, $E(K, N) \approx K$ holds. And as the number of epochs increases to a very large level $(K \gg N^{\frac{b}{a-b}})$, the effective reuse rate eventually converges to a power of $N$. 
% The proof of this theorem can be found in \Cref{sec:proof-one-hot-Ek-power-spectrum}. Beyond this spectrum, we demonstrate that a similar formula holds for other power-law-like spectra. The additional results can be found in \Cref{sec:ont-hot-add}, and its proof can be found in \Cref{sec:proof-one-hot-Ek-log-power-spectrum}.

% \iffalse
\subsection{Results on Logarithmic Power-law Spectrum}\label{sec:log-power-spectrum}

% \haodong{try to understand our results under log spectrum}

Further, we aim to understand under the same Hessian matrix, how the data distribution correlated with $\mP$ and $\Lambda$ affects the effective reusing rate. By changing the spectrum of $\Lambda$, we can also obtain matching upper lower bounds for $\Bar{\cR}^*(K,N)$ and a characterization for $E(K,N)$, which behave differently from the power-spectrum case. Here we present only the latter; the former can be seen in Appendix~\ref{app:log-spectrum}. 

% \vspace{-0.7cm}

\begin{assumption}[Logarithmic Power-Law Spectrum] \label{asp:log-power-decay-prior}
    There exist two constants $a>1,b>0$ such that the data input distribution satisfies that $p_i = c i^{-a} \log^b (i+1)$ and $\Lambda_i = 1/{\log^b (i+1)}$, where $c = \left(\sum_{i=1}^d i^{-a}{\log^b (i+1)}\right)^{-1}$.
\end{assumption}

% \tingkai{add into the first paragraph; by changing some assumptions, derive some results ...}
\begin{theorem}[Multi-Epoch Scaling Under Logarithmic Power-Law Spectrum] \label{thm:one-hot-E_K-log-spectrum}
     Under Assumptions~\ref{asp:parameter-prior}, \Cref{asp:log-power-decay-prior}, and given the data dimension $d = \Omega((KN)^{\frac{1}{a}})$ for a one-pass SGD and a $K$-epoch SGD over $N$ fresh data, it holds that
    \[
    E(K, N) = 
    \begin{cases}
         K(1 + o(1)) & \text{for }  K = o(\log^b N) \\
         \Theta(\log^b N) & \text{for } K = \omega(\log^b N).
    \end{cases}
    \]
\end{theorem}
% In the logarithmic power-law spectrum setting, the dependence of the effective reuse rate on \(K\) closely mirrors both the behavior described in \Cref{thm:convex-E_K} for the strongly convex case and the trend in \Cref{thm:convex-E_K} under the power-law spectrum assumption. In particular, for small \(K\) we still observe an effective-reuse regime ($E_{\eta,\eta}(K, N) \approx K$); as \(K\) grows large, a phase transition occurs and the effective reuse rate converges to an \(N\)–dependent upper bound $\Theta(\log^b N)$

% Under the assumption of a logarithmic power-law spectrum, the trend of the effective reuse rate as a function of \( K \) approximates the phenomena described in \Cref{thm:convex-E_K} in the strongly convex setting and the trend described in \Cref{thm:convex-E_K} under the power-law spectrum assumption. We still observe an effective-reuse regime $(E(K, N) \approx K)$ when \( K \) is relatively small $( K \ll \log^b N)$, and as \( K \) increases, the effective reuse rate undergoes a phase transition, converging to an upper bound determined by \( N \), entering the limited-reuse regime $(E(K, N) =  \Theta(\log^b N))$. Unlike previous theorems, due to the change in problem structure, we observe that the value of \( K \) at which the phase transition occurs, as well as the final upper bound of the effective reuse rate, are dependent on the problem structure.

% \Cref{thm:one-hot-E_K-v2} shows that under a power-law spectrum, the effective reuse rate scales much like in the strongly convex case: when $K \ll \log^b N$, we have $E(K,N)\approx K$, while for $K \gg \log^b N$, it converges to a power of $\log N$, and the exponent of this power is determined by the decay rate of norm decay of the parameter with respect to dimension.

\paragraph{The Saturation Point Varies across Different Problem Setups.} The phase transition point where the effectiveness of data reusing changes from the effective-reuse regime to the limited-reuse regime varies across different problem setups.
% \kaifengclaw{Wording: 'effectively reused/limitedly reused' is awkward; suggest 'from the effective-reuse regime to the limited-reuse regime'.} 
In strongly convex linear regression problems, this phase transition happens when the limit $\lim_{K \to \infty}\frac{K}{\log N}$ changes from $0$ to $\infty$. And in the above power spectrum and log-power spectrum case, the limit turns to be $\lim_{K \to \infty}\frac{K}{N^{b/(a-b)}}$ and $\lim_{K \to \infty}\frac{K}{\log^bN}$.
% \paragraph{General Case: $E(K,N)$ dependent with $N$.} When $b>0$, we show that $E(K, N) \approx K$ when $K$ is relatively small, and saturate to some value dependent of $N$ when $K$ is relatively large. This scaling behavior is identical to that in the strongly-convex case. 

% \iffalse
% \paragraph{Key Proof Insight for Theorem~\ref{thm:one-hot-E_K-v2} and Theorem~\ref{thm:one-hot-E_K-log-spectrum}.} In the small-K regime, we first make approximations based on the following Taylor expansion when $i$ is large:
% \begin{align*}
%     p_i\mu_i^2\left(1-p_i+p_i(1-\eta\mu_i^2)^{2K}\right)^N \approx p_i\mu_i^2\left(1-2\eta Kp_i\mu_i^2)\right)^N
%     \approx p_i\mu_i^2e^{-2\eta NKp_i\mu_i^2}
% \end{align*}
% Then we use Lemma~\ref{lem:E_K} and conclude the proof. In the large-K regime, we have
% \begin{align*}
%     p_i\mu_i^2\left(1-p_i+p_i(1-\eta\mu_i^2)^{2K}\right)^N &\approx p_i\mu_i^2\left(1-p_i\right)^N
% \end{align*}
% We derive matching upper and lower risk bounds for both the ME and OP case, and we finish the proof.
% \fi

% \subsection{Theoretical Insights from The Effective Reuse Rate Formula.} 
% % When $b>0$, we show that $E(K, N) \approx K$ when $K$ is relatively small, and saturate to some value dependent of $N$ when $K$ is relatively large. This scaling behavior is identical to that in the strongly convex case. 
% \haodong{delete the first two points, keep the thrid one}
% \paragraph{Effective-Reuse and Limited-Reuse Regime.} By combining our findings in both the strongly convex case and the non-strongly convex case with a power-law spectrum structure, we observe the general trend of the effective reuse rate:
% \begin{enumerate}
%     \item When the number of epochs $K$ is small and the dataset size $N$ is large, data reuse training enters the effective-reuse regime, where the effective reuse rate is approximately equal to $K$.
%     \item In contrast, when the dataset is reused too many times (i.e., $K$ is large and $N$ is small), the training enters the limited-reuse regime, where the effective reuse rate converges to an upper bound determined by $N$.
%     \item the phase transition point where the effectiveness of data reusing changes from effectively reused to limitedly reused varies across different problem setups. In strongly convex linear regression problems, this phase transition happens when the limit $\lim_{K \to \infty}\frac{K}{\log N}$ changes from $0$ to $\infty$. And in the above power spectrum and log-power spectrum case, the limit turns to be $\lim_{K \to \infty}\frac{K}{N^{b/(a-b)}}$ and $\lim_{K \to \infty}\frac{K}{\log^bN}$.
% \end{enumerate}
% when the number of epochs $K$ is relatively small and the dataset size $N$ is relatively large, the data reusing training enters a effective-reuse regime, where the effective reuse rate approximately equals to $K$
% In contrart, when we repeat a small dataset for too many times, when $K$ is relatively large and the dataset size $N$ is relatively small, the data reusing training enters a limited-reuse regime, where the effective reuse rate converges to an upper bound determined by $N$.

% \paragraph{Larger Datasets Can Be Repeated More.} We have provided a new insight in the calculation of the effective reuse rate in different linear regression problems. Specifically, for larger datasets (as $N$ increases), the permissible $K$ value in the effective-reuse regime also increases. This means that for larger datasets, a higher degree of data reuse is allowed. This observation is further validated in our experiments with large language models in \Cref{sec:experiment-llm}.

% As a side product, our results also theoretically provide counterexamples for the proposed formula for $E_{\eta,\eta'}(K,N)$ in \citet{muennighoff2023scaling}. As shown in \Cref{equ:data-constrained-EK}, which inferred that $E_{\eta,\eta'}(K,N)$ is independent of $N$.

\iffalse
\subsection{Special Case: Isotropic Spectrum}

Finally, we consider a special case, where we assume an isotropic spectrum for $\Lambda$. This situation can be considered an extreme case of data reuse, where the expression for the effective reuse rate is actually independent of the dataset size. The derived expression for $E_{\eta, \eta}(K, N)$ closely resembles the scaling-constrained scaling law form derived by \citet{muennighoff2023scaling}. 

Furthermore, in this setting, when we carefully adjust the learning rate, the training enters the ``complete memorization'' regime, where data reuse becomes entirely ineffective, and the effective reuse rate is the constant $1$.
\begin{assumption}[Isotropic Spectrum] \label{asp:b=0}
    There exists a constant $a>1$ such that the data input distribution satisfies that $p_i = c i^{-a} $ and $\Lambda_i = 1$, where $c = \left(\sum_{i=1}^d\frac{1}{i^{-a}}\right)^{-1}$.
\end{assumption}

\begin{theorem}[Multi-Epoch Scaling Under Isotropic Spectrum] \label{thm:one-hot-E_K-b=0}
    Under Assumptions~\ref{asp:parameter-prior}, \ref{asp:large-dimension} and \ref{asp:b=0}, it holds that $E_{\eta,\eta}(K, N)\eqsim \frac{1-(1-\eta)^{2K}}{2\eta-\eta^2}$ and $E_{1,1}(K, N) = 1$ when $\eta = 1$.
\end{theorem}
% \paragraph{Special Case: $E(K,N)$ independent with $N$.}
The above theorem states that when the data is generated from the standard basis $\{\ve_1,\dots,\ve_d\}$ with a power-law probability matrix, $E_{\eta,\eta}(K, N)\eqsim \frac{1-(1-\eta)^{2K}}{2\eta-\eta^2}$. This formula is similar to the heuristic derivation of the data constrained scaling law in Appendix A of \citep{muennighoff2023scaling}. And hte learning rate $\eta$ corresponds to the proportion of learned information in their derivation. 

Interestingly, the theorem also shows that the case with $\eta = 1$ can be interpreted as a ``complete memorization'' setting where data reuse is useless ($E_{1,1}(K,N)=1$). One could imagine that, in this case, once SGD meets the data $\vx = \ve_i$ for the first time, SGD completely minimizes the loss in this dimension, so that data reuse would not get any benefit in this scenario. 
\fi
% \subsection{Technical Challenges and Key Proof Insights}
% The proof of logarithmic power-law spectrum case follows a similar pattern to that of power-law spectrum case. So we give the proof sketch of power-law spectrum case, and only list the lemmas of logarithmic power-law spectrum case in the following.
% \subsubsection{Derivation of the Optimal Learning Rate}
% Given the precise formula in \Cref{thm:one-hot-loss-estimate}, we can see that the excess risk strictly decreases with $\eta$ when $\eta \in [0, \frac{1}{D^2}]$. Nevertheless, the dependency between excess risk and learning rate still remains elusive when $\eta \in [\frac{1}{D^2}, \frac{2}{D^2}]$ due to distinct variation trends of different terms as $\eta$ increases, which makes the analysis for optimal learning rate extermely challenging. To tackle with the issue, here we give a characterization of a learning rate set that achieves a near-optimal excess risk.
% \begin{lemma} \label{lem:one-hot-power-small-K-opt-lr}
%     Suppose the assumptions in \Cref{thm:one-hot-E_K-v2} hold. When $K = o(N^{\frac{b}{a-b}})$, for all $\eta \in \left[2-o(1), 2-\omega(\frac{\log KN}{KN})\right]$, we have $\Bar{\cR}(K, N;\eta)=\cR^*(K, N)(1+o(1))$.
% \end{lemma}
% \begin{lemma} \label{lem:one-hot-power-large-K-opt-lr}
%     Suppose the assumptions in \Cref{thm:one-hot-E_K-v2} hold. When $K = \omega(N^{\frac{b}{a-b}})$, for all $\eta \in \left[1, 2\right]$ we have $\cR^*(K, N)\eqsim...$\tingkai{we only need this.}
% \end{lemma}

%  We present the lemmas for logarithmic spectrum case below.
% \begin{lemma} \label{lem:one-hot-log-power-small-K-opt-lr}
%     Suppose the assumptions in \Cref{thm:one-hot-E_K-v2} hold. When $K = o(\log^b N)$, for all $\eta \in \left[2\log^b2-o(1), 2\log^b2-?\right]$, we have $\Bar{\cR}(K, N;\eta)=\cR^*(K, N)(1+o(1))$.
% \end{lemma}
% \begin{lemma} \label{lem:one-hot-log-power-large-K-opt-lr}
%     Suppose the assumptions in \Cref{thm:one-hot-E_K-v2} hold. When $K = \omega(\log^b N)$, for all $\eta \in \left[\log^b2, 2\log^b2-?\right]$ we have $\Bar{\cR}(K, N;\eta)=\cR^*(K, N)(1+o(1))$.
% \end{lemma}

% \subsubsection{Derivation of $E(K, N)$}
% \fi